\باب{کثیر  استعمال حدود}\شناخت{ضمیمہ_و}

یہ ضمیمہ  حصہ \حوالہء{11.1}  کے مسئلہ \حوالہء{5}   کے حد  \اردوعدد{4} تا \اردوعدد{6} کی تصدیق کرتا ہے۔

\ابتدا{ثبوت}\موٹا{حد \اردوعدد{4}:}\quad \موٹا{اگر \عددی{|x|<1} ہو تب \عددی{\lim_{x^n}x^n=0} ہو گا۔}\quad 
ہم نے دکھانا ہو گا کہ   ہر \عددی{\epsilon>0} کا   مطابقتی    اتنا بڑا عدد صحیح \عددی{N} موجود ہے  جس سے بڑے تمام  اعداد \عددی{n} کے
 لئے \عددی{|x^n|<\epsilon} ہو گا۔ چونکہ  \عددی{|x|<1} کی صورت میں \عددی{\epsilon^{1/n}\to 1} ہے ، ایسا عدد صحیح موجود ہو گا جس کے لئے \عددی{\epsilon^{1/N}>|x|} ہو۔دوسرے الفاظ میں درج ذیل ہو گا۔
\begin{align}\label{مساوات_ضمیمہ_ج_ایک}
|x^N|=|x|^N<\epsilon
\end{align}
یہ وہ عدد صحیح ہے جس کی ہمیں تلاش ہے چونکہ \عددی{|x|<1} کی صورت میں درج ذیل ہو گا۔
\begin{align}\label{مساوات_ضمیمہ_ج_دو}
|x^n|<|x^N|\quad \text{\RL{تمام $n>N$ کے لئے}}
\end{align}
مساوات \حوالہ{مساوات_ضمیمہ_ج_ایک} اور مساوات \حوالہ{مساوات_ضمیمہ_ج_دو} ملا کر ، تمام \عددی{n>N} کے 
لئے   \عددی{|x^n|<\epsilon} حاصل ہو گا۔یوں ثبوت مکمل ہوا۔
\انتہا{ثبوت}

\ابتدا{ثبوت} \موٹا{حد \اردوعدد{5}:}\quad
 \موٹا{کسی بھی عدد \عددی{x} کے لئے \عددی{\lim_{n\to \infty}(1+\tfrac{x}{n})^n=e^x} ہو گا۔} \\
 درج ذیل فرض کرتے ہوئے
\[a_n=\left(1+\frac{x}{n}\right)^n\]
درج ذیل لکھا جا سکتا ہے۔
\[\ln a_n=\ln \left(1+\frac{x}{n}\right)^n=n\ln\left(1+\frac{x}{n}\right)\]

\ترچھا{ال ہوس پٹل}  کا قاعدہ  استعمال کرتے ہوئے،  ہم \عددی{n} کے لحاظ سے تفرق لے کر حد تلاش کرتے ہیں۔
\begin{align*}
\lim_{n\to\infty} n\ln \left(1+\frac{x}{n}\right) &=\lim_{n\to\infty}\frac{\ln (1+x\!/\!n)}{1\!/\!n}\\
&=\lim_{n\to\infty}\frac{\left(\frac{1}{1+x\!/\!n}\right)\cdot\left(-\frac{x}{n^2}\right)}{-1\!/\!n^2}=\lim_{n\to\infty}\frac{x}{1+x\!/\!n}=x
\end{align*}
یوں درج ذیل ہو گا۔
\[\ln a_n\to x\]
حصہ \حوالہء{11.1} کا مسئلہ \حوالہء{4} بروئے کار لاتے ہوئے اور \عددی{f(x)=e^x} لے کر درج ذیل اخذ ہو گا۔
\[\left(1+\frac{x}{n}\right)^n=a_n=e^{\ln a_n}\to e^x\]
\انتہا{ثبوت}

\ابتدا{ثبوت}\موٹا{کسی بھی عدد \عددی{x} کے لئے \عددی{\lim_{n\to\infty}\tfrac{x^n}{n!}=0} ہو گا۔}\\
درج ذیل جانتے ہوئے
\[-\frac{|x|^n}{n!}\le \frac{x^n}{n!}\le \frac{|x|^n}{n!}\]
ہمیں صرف اتنا دکھانا ہے کہ \عددی{|x|^n\!/\!n!\to 0} ہو گا۔ اس کے بعد ہم   مسئلہ بیچ برائے  ترتیبات    (حصہ \حوالہء{11.1} میں    مسئلہ  \حوالہء{2})  استعمال کر کے  \عددی{x^n\!/\!n!=0} اخذ کر سکتے ہیں۔

یہ دکھانے کی خاطر کہ \عددی{|x|^n\!/\!n!\to 0} ہے ، سب سے پہلے  ایسا عدد صحیح \عددی{M>|x|}  منتخب کرنا ہو گا   جو \عددی{(|x|/M)<1} دے۔  یوں حد \اردوعدد{4} جو ہم ثابت کر چکے کے تحت \عددی{(|x|/M)^n\to 0} ہو گا۔لہٰذا   \عددی{n>M} قیمتوں تک اپنی توجہ محدود رکھ کر ہم   درج ذیل لکھ سکتے ہیں۔
\begin{align*}
\frac{|x|^n}{n!}&=\frac{|x|^n}{1\cdot 2\cdots M\cdot \underbrace{(M+1)(M+2)\cdot n}_{\text{\RL{$(n-M)$ جزو ضربی}}}}\\
&\le \frac{|x|^n}{M!M^{n-M}}=\frac{|x|^n M^M}{M!M^n}=\frac{M^M}{M!}\left(\frac{|x|}{M}\right)^n
\end{align*}
یوں درج ذیل ہو گا۔
\[0\le \frac{|x|^n}{n!}\le \frac{M^M}{M!}\left(\frac{|x|}{M}\right)^n\]
اب جیسے جیسے \عددی{n} کی قیمت بڑھتی ہے مستقل \عددی{M^M/M!} کی قیمت تبدیل نہیں ہوتی۔ یوں مسئلہ بیچ کے تحت \عددی{(|x|/M)^n\to 0} کی وجہ سے \عددی{|x|^n/n!\to 0} ہو گا۔
\انتہا{ثبوت}
